\section{Count Square Submatrices with All Ones}
\label{sec:dp-count-square-submatrices}

\problemstatement{Given a binary matrix $\textit{mat}$, count the total
number of square submatrices that contain only $1$s (squares of every
size, at every position).}

\begin{patternbox}
\emph{``Largest/count of square submatrices of all ones, anchored at
each cell''} is solved by a recurrence that looks at exactly three
neighbors: a square of side $s$ ending at $(i,j)$ as its bottom-right
corner exists only if squares of side $s-1$ already exist ending at
\textbf{all three} of the cell above, the cell to the left, and the
cell diagonally above-left.
\end{patternbox}

\subsection*{Deriving the recurrence}

Let $\textit{dp}[i][j]$ be the side length of the \textbf{largest}
all-ones square whose bottom-right corner is exactly $(i,j)$ (and
$0$ if $\textit{mat}[i][j]=0$). If $\textit{mat}[i][j]=1$, extending a
square one row down and one column right requires the square
\emph{ending} at each of the three neighboring corners to already
support at least side length $\textit{dp}[i][j]-1$:
\[
  \textit{dp}[i][j] = \begin{cases}
  0 & \textit{mat}[i][j]=0, \\
  1+\min\big(\textit{dp}[i{-}1][j],\ \textit{dp}[i][j{-}1],\
  \textit{dp}[i{-}1][j{-}1]\big) & \textit{mat}[i][j]=1.
  \end{cases}
\]

\begin{ideabox}[Why the Minimum, and Why Summing Gives the Count]
The new square's side is capped by whichever of the three neighbors
supports the \emph{smallest} square -- a single weak corner limits how
far the square can extend, exactly like the weakest link in a chain.
Separately, a cell with $\textit{dp}[i][j]=k$ doesn't just mean
\emph{one} square ends there -- it means $k$ different squares end
there: one each of side $1, 2, \ldots, k$ (every smaller square nested
inside the largest one is also a valid, distinct all-ones square with
its bottom-right corner at $(i,j)$). So the \textbf{count} of all
square submatrices is simply $\sum_{i,j} \textit{dp}[i][j]$.
\end{ideabox}

\subsection*{C++ implementation}

\begin{lstlisting}
#include <bits/stdc++.h>
using namespace std;

int countSquareSubmatrices(vector<vector<int>>& mat) {
    int n = mat.size(), m = mat[0].size();
    vector<vector<int>> dp(n, vector<int>(m, 0));
    int totalCount = 0;

    for (int i = 0; i < n; i++) {
        for (int j = 0; j < m; j++) {
            if (mat[i][j] == 0) {
                dp[i][j] = 0;
            } else if (i == 0 || j == 0) {
                dp[i][j] = 1; // first row/column: at most a 1x1 square
            } else {
                dp[i][j] = 1 + min({dp[i - 1][j], dp[i][j - 1], dp[i - 1][j - 1]});
            }
            totalCount += dp[i][j];
        }
    }
    return totalCount;
}

int main() {
    vector<vector<int>> mat = {
        {0, 1, 1, 1},
        {1, 1, 1, 1},
        {0, 1, 1, 1}
    };
    cout << countSquareSubmatrices(mat) << endl; // 15
    return 0;
}
\end{lstlisting}

\subsection*{Dry run}

\dryrun{Take the $3\times4$ matrix above.}

\begin{center}
\begin{tabular}{c|c|c|c|c}
 & 0 & 1 & 2 & 3 \\ \hline
0 & 0 & 1 & 1 & 1 \\ \hline
1 & 1 & 1 & 2 & 2 \\ \hline
2 & 0 & 1 & 2 & 3
\end{tabular}
\end{center}

$\textit{dp}[1][2]=2$: $\textit{mat}[1][2]=1$, and
$\min(\textit{dp}[0][2],\textit{dp}[1][1],\textit{dp}[0][1])=
\min(1,1,1)=1$, so $\textit{dp}[1][2]=2$ (a $2\times2$ all-ones square
exists with corners at rows $0$--$1$, columns $1$--$2$).
$\textit{dp}[2][3]=3$: neighbors $\textit{dp}[1][3]=2$,
$\textit{dp}[2][2]=2$, $\textit{dp}[1][2]=2$, minimum $2$, giving
$\textit{dp}[2][3]=3$ (a full $3\times3$ all-ones square, columns
$1$--$3$, all three rows). Summing every cell:
$0{+}1{+}1{+}1{+}1{+}1{+}2{+}2{+}0{+}1{+}2{+}3 = 15$.

\begin{complexitybox}
\textbf{Time Complexity.} $O(n \cdot m)$ -- one constant-time
comparison per cell.
\textbf{Space Complexity.} $O(n \cdot m)$ for the \textit{dp} table
(reducible to $O(m)$ with a rolling pair of rows, since each cell only
reads the row directly above and the current row).
\end{complexitybox}

\begin{takeawaybox}
A surprisingly large class of ``2D shape ending here'' problems reduce
to a three-neighbor minimum (or maximum) plus one: the moment a shape
can only grow when \emph{all} of its supporting smaller shapes already
exist, the recurrence is this same one-line comparison, with no search
over a range needed at all -- the simplest recurrence shape in this
entire book.
\end{takeawaybox}
